% !TEX root = ../notes_bizon.tex
% Chapter 1: Quickstart -- from your laptop to a shell on the Bizon
%%%%%%%%%%%%%%%%%%%%%%%%%%%%%%%%%%%%%%%%%%%%%%%%%%%%%%%%%%%%%%%%%%%%%%
%
% Access path: Sophos SSL VPN into the DHBW network, then ssh with the
% DHBW account, as in the DHBW guide (bizon.bib). Students log in as
% <username>@ab.ba.ba-ravensburg.de. The server address never appears in
% the book: it ships in the setup package (student_setup/, gitignored),
% handed out by mail as for the DGX Spark.
% TODO after bizon_recon.sh: confirm the ssh -l form and the GPU count.

\index{quickstart}\index{Bizon}
\chapter{Quickstart.}
\printchapterversion{01_quickstart}

\section{Where the buffalo roam}

\index{account}\newthought{The Bizon runs on your DHBW account.} Log in
with the username and password of the lab PCs. On the server, the
username carries the domain: \texttt{<username>@ab.ba.ba-ravensburg.de}
\cite{dhbw2021bizon}. The Bizon's address comes by mail, as a setup package.
\marginnote{\newthought{Scan the code} and send a mail stating your name,
course, and what you want to run.\par\medskip
\null\hfill\qrcode[height=1.0cm]{mailto:schutera@dhbw-ravensburg.de?subject=Bizon\%20Setup\&body=Name\%3A\%20\%0ACourse\%3A\%20\%0AProject\%3A\%20}}

\index{shared infrastructure}\newthought{Your account is one of many} on
a shared machine. Store nothing sensitive on it, leave other users' files
and processes alone, and sandbox your stuff.
Protect yourself from others, and protect others from yourself.
\marginnote{\newthought{Remember?} Chapter 1 of
Full Stack Handwerkszeug, Development and Production, covers this split,
along with SSH and Linux servers.
\url{https://material.schutera.com/lecturenotes/notes_missingsemester/notes_missingsemester.pdf}}%
The Bizon is used most efficiently when you spend as little time as possible to deploy your jobs and leave,
not for interactive development. If you need to develop interactively, please do so on your own machine.

\marginnote{\newthought{Scan the code} for the VPN user portal.\par\medskip
\null\hfill\qrcode[height=1.0cm]{https://vpn-rv.dhbw-ravensburg.de}}
\index{VPN}\newthought{The Sophos VPN takes you into the DHBW network.}
On campus you are already inside. Off campus, open
\url{https://vpn-rv.dhbw-ravensburg.de}, log in to the user portal with
your DHBW credentials, open the page \texttt{VPN}, and download the
client and the configuration. Then connect with your DHBW credentials:
\begin{itemize}[nosep]
  \item Windows: Sophos SSL VPN Client, tray icon, \texttt{Connect}.
  \item macOS: \texttt{.ovpn} into Tunnelblick \cite{tunnelblick2026}, \texttt{Connect}.
  \item Linux: \texttt{sudo openvpn -{}-config <file>.ovpn}.
\end{itemize}

\index{SSH}\newthought{Then log in} with the address from the package:
\begin{verbatim}
ssh -l <username>@ab.ba.ba-ravensburg.de <bizon-ip>
\end{verbatim}

\index{nvidia-smi}\newthought{Before you start a job,} run
\texttt{nvidia-smi} to see who is on which GPU. Pick a free one and pin
your job to it, for example GPU 3:
\begin{verbatim}
CUDA_VISIBLE_DEVICES=3 python train.py
\end{verbatim}
